%! Orthogonality of the Four Subspaces — Unit 4, Lesson 4.1 — Homework
\documentclass[10pt]{article}
\usepackage{linalg-article}
\usepackage{linalg-boxes}
\newcommand{\TallMath}[1]{$\displaystyle #1\rule[-1.4em]{0pt}{3.2em}$}
\newcommand{\vv}{\mathbf{v}}
\newcommand{\ww}{\mathbf{w}}
\newcommand{\xx}{\mathbf{x}}
\newcommand{\yy}{\mathbf{y}}
\newcommand{\bb}{\mathbf{b}}
\newcommand{\pp}{\mathbf{p}}
\newcommand{\avec}{\mathbf{a}}
\newcommand{\zero}{\mathbf{0}}
\newcommand{\T}{\mathsf{T}}

\begin{document}

\pageheader{Unit 4, Lesson 4.1}{Homework}
\namedateperiod

\vspace{0.12in}

\begin{notesbox}{Practice}
Two subspaces are \textbf{orthogonal} when every vector in one is perpendicular to every vector in the
other; they are \textbf{complements} when they are also of dimensions that fill the space. Show your
dot products.

\begin{enumerate}[leftmargin=1.9em, itemsep=11pt]
  \item \textbf{Perpendicular to a whole subspace.} A plane $V\subset\mathbb{R}^3$ is spanned by
        $\begin{bmatrix} 1\\0\\2 \end{bmatrix}$ and $\begin{bmatrix} 0\\1\\1 \end{bmatrix}$. Show that
        $\ww=\begin{bmatrix} 2\\1\\-1 \end{bmatrix}$ is perpendicular to the \emph{entire} plane by
        dotting it with \emph{both} spanning vectors. \writelines{2}
  \item \textbf{Both rooms of one matrix.} Let
        $C=\begin{bmatrix} 1 & 0 & 1 \\ 0 & 1 & 1 \\ 1 & 1 & 2 \end{bmatrix}$ (rank $r=2$), with
        nullspace direction $\xx=\begin{bmatrix} -1\\-1\\1 \end{bmatrix}$ and left-nullspace direction
        $\yy=\begin{bmatrix} 1\\1\\-1 \end{bmatrix}$.
    \begin{enumerate}[label=(\alph*), leftmargin=1.8em, itemsep=5pt]
      \item Dot $\xx$ with each \emph{row} of $C$. What pairing does this confirm? \vspace{0.9cm}
      \item Dot $\yy$ with each \emph{column} of $C$. What pairing does this confirm? \vspace{0.9cm}
    \end{enumerate}
  \item \textbf{Count \& name.} An $m\times n$ matrix has $m=5$, $n=3$, rank $r=3$. Give the dimension
        of all four subspaces. Which pair are complements in $\mathbb{R}^3$? in $\mathbb{R}^5$?
        \writelines{2}
  \item \textbf{Perpendicular vs.\ complement.} Explain the difference between two subspaces being
        \emph{perpendicular} and being \emph{orthogonal complements}. Give a $\mathbb{R}^3$ example that
        is perpendicular but \emph{not} a complement. \writelines{2}
\end{enumerate}
\end{notesbox}

\begin{extensionbox}
\textbf{The complement of a line.} Let $L$ be the line in $\mathbb{R}^3$ through
$\avec=\begin{bmatrix} 1\\2\\2 \end{bmatrix}$.
\begin{enumerate}[label=(\alph*), leftmargin=1.8em, itemsep=5pt]
  \item $L^{\perp}$ is all $\ww$ with $\avec\cdot\ww=0$. Write that single equation in $w_1,w_2,w_3$, and
        give \emph{two} independent vectors that satisfy it.
  \item What are $\dim L$ and $\dim L^{\perp}$? Do they fill $\mathbb{R}^3$? Explain why $L^{\perp}$ is the
        \emph{nullspace} of the $1\times 3$ matrix $A=\begin{bsmallmatrix} 1 & 2 & 2 \end{bsmallmatrix}$.
\end{enumerate}
\writelines{3}
\end{extensionbox}

\begin{spiralbox}
\textbf{Coming next --- Lesson 4.2: Projections.} Today the error of a best guess will always land in an
orthogonal complement. Next we \emph{use} that: to find the closest point of a line or subspace to a
target $\bb$, we force the error $\bb-\pp$ into the complement (dot $=0$) and solve --- the projection.
\end{spiralbox}

\end{document}
